% Nouveau EGA V (cote 25): archive cote 25, site page 109 (PDF page 110).
% typescript. Draft compared: none.
% AI-transcribed by Claude Opus 5.5 (claude-opus-5-5, adaptive thinking, effort high), Message Batches API, batch ega25b (msgbatch_01A7n5YuXubUPNK1zsepi8RG), 2026-11-13; checked against both supplied scan parts; see CHECK_DECISIONS.tsv.
\scanpage{25/109}
\mspage{12}

en vue du par.25 sur le résult \unsure{ant|aht} et le discriminant des formes, sera la
déter\var{\over{m}{m}}{m}{writing over}{ega25_25_109-v0001}ination des diviseurs relatifs $\hand{\geqslant} 0$ sur un fibré multiprojectif.
\var{\inserted[typed]{2.8.}}{2.8.}{insertion}{ega25_25_109-v0002}
Si $\underline{D}$ est un diviseur relatif sur P, on définit son multidegré comme étan
% NOTE: the right edge of the page is cut off; "étan[t]" and "ramené[e]" lose their last letters
celui de $\underline{O}_P(D)$. Comme plus haut, la détermination de $\mathrm{Div}^{+}(P/S)$ est ramené
à celle de $\mathrm{Div}^{\underline{n}}(P/S)$, pour un multidegré $\underline{n}\hand{\in}\underline{\unsure{N|M}}^{I}$ donné, qui fournit un
isomorphisme
\[ \underline{L} =\underline{O}_X(D) \hand{\simeq} \underline{M} \otimes_{\underline{O}_S} \underline{O}_P(\underline{n}) \ . \]
Or on obtient, grâce au Chap II par.2 (\var{\xout{ou à Künneth}}{}{typewriter strike-out}{ega25_25_109-v0003}
\[ (xx) \qquad f_{\hand{*}}(\underline{D}) = \underline{M} \otimes f_{\hand{*}}(\underline{O}_P(\underline{n})) = \underline{M} \otimes \hand{\bigotimes}_i \underline{\mathrm{Sym}}^{n_i}(\underline{E}_i) \quad . \]
Proéédant maintenaht comme pour 2.2. on trouve la

\uline{Proposition} 2.9. Avec les notations précédentes, à reppaler, on a\var{\xout{\illegible{1}}}{}{typewriter strike-out}{ega25_25_109-v0004} une
corr biun can entre l'ensemble $\mathrm{Div}^{\underline{n}}(P/S)$ des diviseurs relatifs de
multidegré $\underline{n}$ sur P, et l'ensemble des Modules inversibles quotients $\underline{M}$
de $\hand{\bigotimes}_i \underline{\mathrm{Sym}}^{n_i}(\underline{E}_i)^{\hand{\vee}}$ (ou ce qui revient au même ...) . Si D et $\underline{M}$ se corres-%
pondent, $\underline{O}_P(D)$ est can isomorphe à $\underline{M}\otimes\underline{O}_P(\underline{n})$ , et $s_D$ s'identifie alors
à ce qu'on devine, compte tenu de (xx).

\uline{Corollaire} 2.10. Le sous-foncteur $\mathrm{Div}^{\underline{n}}_{P/S}$ de $\mathrm{Div}^{+}_{P/S}$ correspondant aux
diviseurs relatifs de multidegré $\underline{n}$ est can représentable par le fibré
projectif $P(\hand{\bigotimes}_i \underline{\mathrm{Sym}}^{n_i}(\underline{E}_i)^{\hand{\vee}})$, et $\underline{\mathrm{Div}}^{+}_{P/S}$ est can repr par le préschéma
somme de ces de\unsure{rni|mmi}ers, pour $\underline{n}\hand{\in}\underline{N}^{I}$ .

2.11. La détermination précédente très simple des $\mathrm{Div}_{P/S}$ est due à la
structure très simple de Pic(P/S) (en fait, à la structure``discrète'' du
préschéma de Picard $\underline{\mathrm{Pic}}_{P/S}$ ....). Nous pouvons abstraire le raisonnement
fait, qui revient essentiellement à établir une représentabilité ``relative''
(par rapport à $\underline{\mathrm{Pic}}$ ). P$_o$ur ceci, donnons-nous un morphisme
\[ f\colon X \hand{\to} S \]
propre et plat de présentation finie, \uline{et} un Module inversible $\underline{L}$ sur X.

\mspage{109}
